\section{The \lib{} library}
\label{sec:library}

\lib{} is a Python package whose base install depends only on numpy, pandas and scipy. It
has three layers. \emph{Skills} hold the financial knowledge a model needs and the
algorithms that apply it. \emph{Checks} are executable tests for the failures that the
knowledge is meant to prevent. \emph{Self-improvement} methods turn check results into
better programs (\S\ref{sec:methods}). An agent reaches all three through the same
interfaces (\S\ref{sec:interfaces}; Figure~\ref{fig:overview}).

\subsection{Skills: documents, algorithms and rule cards}
\label{sec:skills}

\paragraph{Skill documents.}
The library ships \NSkills{} skills in \NPlugins{} plugins that follow the Agent Skills
specification \citep{agentskills}, with \NReferenceFiles{} reference files and
\NScripts{} runnable scripts. A skill is not a tutorial. It records the defaults, licence
changes and silent failure modes that a model's training prior gets wrong. Each claim is
dated and marked by how it was verified: at a primary source, secondhand, silently wrong,
or dead. The descriptions of all \NSkills{} skills take about \NListingTokens{} tokens of an
agent's skill listing, so the skills are split into plugins and a user installs only the
ones a task needs.
The plugin marketplace also lists \NFederatedPacks{} federated packs from other repositories
(\NFederatedSkills{} skills), off by default and not verified here. \NLensPacks{} of them are
investor templates: \NLensSkills{} skills built on named investors' approaches, such as
Buffett's, Soros's or Duan Yongping's. The \texttt{kol-credibility-registry} skill adds track
records for \NKOLProfiles{} social-media authors.

\paragraph{Algorithms.}
A catalog of \NAlgorithms{} algorithm entries covers \NAlgorithmTasks{} task types,
including portfolio construction, trading signals, volatility, risk, forecasting, pricing and
execution. \NNative{} are implemented in the library. The other \NAdapters{} wrap
\NOptionalLibraries{} optional libraries; when a library is absent, its entries are reported
as missing dependencies rather than replaced silently. Two functions turn one algorithm into
the object a research program needs:
\begin{lstlisting}[language=Python]
w = positions("macd", raw, factors,
              parameters={"fast": 8})
v = rolling("ewma_volatility", raw,
            factors, window=40)
\end{lstlisting}
\texttt{positions()} takes raw quotes and split factors, applies the factors once, and
returns positions decided on the previous close. \texttt{rolling()} does the same for
volatility, risk and forecast algorithms, with one lagged value per date and asset.
\texttt{adjust\_prices()} marks its output as adjusted, and \texttt{positions()} refuses to
adjust such a frame again, because passing adjusted prices together with their factors
applies every split twice.

\paragraph{Rule cards.}
Between the documents and the code sit \NCards{} rule cards in \NCardFamilies{} families:
timing, corporate actions and data hygiene (Table~\ref{tab:cards}). A card holds:
\begin{itemize}
\item a few sentences of rule text;
\item the skill passage the rule comes from;
\item the checks that verify it;
\item the task traits that select it;
\item where one exists, the library call that follows the rule.
\end{itemize}
A task declares its traits, such as \emph{positions} or \emph{split factors}, and receives
only the cards those traits select. A run without library access receives the same cards
without the library calls.

\begin{table}[t]
\centering
\scriptsize
% Generated by make_numbers.py from fin_skills.api.rules; do not edit by hand.
\setlength{\tabcolsep}{4pt}
\begin{tabular}{@{}l>{\raggedright\arraybackslash}p{3.0cm}c@{}}
\toprule
Card & Verified by & Call \\
\midrule
\multicolumn{3}{@{}l}{\emph{timing}} \\
\quad \texttt{execution\_lag} & \texttt{execution\_lag} & yes \\
\quad \texttt{feature\_causality} & \texttt{assert\_causal} & -- \\
\multicolumn{3}{@{}l}{\emph{corporate actions}} \\
\quad \texttt{split\_once} & \texttt{split\_representation} \newline \texttt{adjustment\_check} & yes \\
\quad \texttt{same\_convention} & \texttt{split\_representation} & -- \\
\multicolumn{3}{@{}l}{\emph{data hygiene}} \\
\quad \texttt{warmup} & \texttt{finite\_output} & -- \\
\quad \texttt{first\_return} & \texttt{finite\_output} & -- \\
\bottomrule
\end{tabular}

\caption{Rule cards by family, the checks that verify each one, and whether the card names a
library call that follows the rule.}
\label{tab:cards}
\end{table}

\subsection{Checks}
\label{sec:checks}

The library has \NGuards{} guards behind one interface: \texttt{run(**inputs)} returns an
outcome with \texttt{passed}, a message and warnings. A failing guard is a result, not an
error. Malformed arguments are reported as tool errors that name the field, so a model can
correct its call. \NGuardsExported{} guards take JSON inputs and are exported as tools. The
other \NGuardsPythonOnly{} need a live Python object, such as the position function under
test. A tool schema cannot carry one, and a stand-in built by the model would not be the
caller's function, so these guards are available only in Python.

For generated programs, the library provides property checks that need no expected output.
Three are used throughout this paper:
\begin{itemize}
\item \texttt{finite\_output}: the result has no missing or infinite values.
\item \texttt{split\_representation}: the program is re-run on the same economic prices,
given as adjusted closes with unit factors, and its output must not change (absolute
tolerance $10^{-8}$). A program that applies the factors twice, or mixes raw and adjusted
prices, fails.
\item \texttt{execution\_lag}: the prices from row $k$ onward are perturbed, and no position
up to row $k$ may move (Figure~\ref{fig:episode}c). A rule that trades on the close it reads
fails this check. It still passes a plain causality test, which lets the value on row $t$ use
row $t$.
\end{itemize}
Each check carries a \emph{cause}: one sentence saying what the check observed, for
instance that a position moved when only prices on or after its own row changed.
\texttt{explain()} pairs each cause with the cards that address it. We treat the cause as
part of the check and the card as part of the skills; \S\ref{sec:arms} separates the two.

\subsection{Agent interfaces}
\label{sec:interfaces}

One registry yields \NTools{} tools under the same names in three formats: a Model Context
Protocol server over stdio \citep{mcp2026spec}, OpenAI tool schemas and Anthropic tool
schemas. The tools read and search skills; list, recommend and run algorithms; and run the
exported guards. Python code imports the same functions directly.

\texttt{agents.Team} turns one user model, a callable \texttt{generate(messages) -> str}, into
a writer and one reviewer per rule family. Code checks decide acceptance, and no role can
accept a program the checks reject. The self-improvement methods take the same
\texttt{generate} callable, so a user connects one model to every part of the library.
